% GENERATED by tools/coverage.py from the lesson plans -- do not edit by hand.
\subsection{Objectives per goal and week (generated)}
Each number is how many of that week's \enquote{The student can\ldots} objectives
serve the goal. Blank = not a focus that week (it may still be practiced).

{\small\setlength{\tabcolsep}{3pt}
\begin{longtable}{@{}lcccccccccccccc@{}}
\toprule
Goal & 1 & 2 & 3 & 4 & 5 & 6 & 7 & 8 & 9 & 10 & 11 & 12 & B1 & B2\\ \midrule\endhead
L2.G1 & 2 &  &  &  &  &  &  &  &  &  &  &  &  & 1\\
L2.G2 & 1 &  & 1 & 2 &  &  &  &  & 1 & 1 &  &  & 1 & \\
L2.G3 &  & 3 &  &  &  &  &  & 1 & 1 & 1 &  & 1 &  & 1\\
L2.G4 &  &  & 2 &  & 2 &  & 1 &  & 1 &  & 1 &  &  & 1\\
L2.G5 &  &  &  &  &  &  & 2 &  &  &  & 1 & 1 &  & \\
L2.G6 & 1 &  &  &  & 1 & 4 & 1 & 4 &  &  & 1 & 1 &  & \\
L2.G7 &  & 1 &  & 1 &  &  &  &  & 2 & 1 &  & 1 &  & 2\\
L2.G8 &  &  &  &  &  &  &  &  & 1 & 3 & 1 & 2 & 2 & \\
L2.G9 &  &  & 1 & 1 & 1 &  & 1 & 1 & 1 &  & 1 & 1 & 1 & 1\\
\bottomrule\end{longtable}}

\subsection{All lesson objectives (generated)}
\begin{longtable}{@{}p{1.1cm}p{0.78\linewidth}p{1.6cm}@{}}
\toprule Week & Objective & Goals\\ \midrule\endhead
\multicolumn{3}{@{}l}{\textbf{1: Welcome Back, Python}}\\
1.1 & The student can sort values into the four types \texttt{int}, \texttt{float}, \texttt{str} and \texttt{bool} and check a guess with \texttt{type()}. & L2.G1\\
1.2 & The student can build a message with an f-string and turn typed text into a number with \texttt{int()}. & L2.G1\\
1.3 & The student can combine conditions with \texttt{and}, \texttt{or} and \texttt{not} inside \texttt{if}/\texttt{elif}/\texttt{else}. & L2.G2\\
1.4 & The student can read a \texttt{TypeError} from the bottom up and explain why \texttt{"5" + 5} fails. & L2.G6\\
\midrule
\multicolumn{3}{@{}l}{\textbf{2: Functions That Give Back}}\\
2.1 & The student can explain the difference between printing an answer and returning it, and use a returned value in a new calculation. & L2.G3\\
2.2 & The student can write a function with parameters that returns an answer, such as \texttt{damage(attack, shield)}. & L2.G3\\
2.3 & The student can call one drawing function several times with different arguments to build a picture. & L2.G7\\
2.4 & The student can write a docstring that says what a function gives back, and say why a variable made inside a function is not visible outside it. & L2.G3\\
\midrule
\multicolumn{3}{@{}l}{\textbf{3: Lists: Collections of Things}}\\
3.1 & The student can make a list, add to it with \texttt{append} and count it with \texttt{len}. & L2.G4\\
3.2 & The student can get an item by its index, starting from 0, and explain why \texttt{items[len(items)]} gives an \texttt{IndexError}. & L2.G4\\
3.3 & The student can loop over a list with \texttt{for item in items} and with \texttt{for i in range(len(items))}. & L2.G2\\
3.4 & The student can choose a color-blind-safe palette and label a picture so it never relies on color alone. & L2.G9\\
\midrule
\multicolumn{3}{@{}l}{\textbf{4: While Loops and Game Loops}}\\
4.1 & The student can write a \texttt{while} loop that stops because a variable changes, and explain why a loop that never changes runs forever. & L2.G2\\
4.2 & The student can use \texttt{while True} with \texttt{break} or \texttt{return} to keep asking until an answer is valid. & L2.G2\\
4.3 & The student can use a loop inside a loop to visit every pixel of a grid made from a list of strings. & L2.G7\\
4.4 & The student can say what makes a game fair and name one trick a game should not play on its players. & L2.G9\\
\midrule
\multicolumn{3}{@{}l}{\textbf{5: Dictionaries and Secret Codes}}\\
5.1 & The student can make a dictionary, read, change and add keys, and loop over it with \texttt{.items()}. & L2.G4\\
5.2 & The student can encode and decode a message with a Caesar shift, on paper and in Python. & L2.G4\\
5.3 & The student can explain what a \texttt{KeyError} means and avoid it with \texttt{if key in dictionary}. & L2.G6\\
5.4 & The student can explain why passwords and other secrets need protecting, and never type a real password into our programs. & L2.G9\\
\midrule
\multicolumn{3}{@{}l}{\textbf{6: Review and Checkpoint}}\\
6.1 & The student can read a traceback from the bottom up and name the error type and the line where it happened. & L2.G6\\
6.2 & The student can use \texttt{print()} lines and a trace table to find where a program's values go wrong. & L2.G6\\
6.3 & The student can explain code line by line to a partner or a rubber duck to find a bug that does not crash. & L2.G6\\
6.4 & The student can fix bugs that involve types, list indexes, loops, return values and dictionary keys from weeks 1--5. & L2.G6\\
\midrule
\multicolumn{3}{@{}l}{\textbf{7: Files: Saving Your Stuff}}\\
7.1 & The student can save data to a text file with \texttt{open} and \texttt{with}, choosing \texttt{"a"} to add and \texttt{"w"} to replace. & L2.G5\\
7.2 & The student can read a file line by line and clean each line with \texttt{strip()} and \texttt{split()} into a list. & L2.G4, L2.G5\\
7.3 & The student can handle a missing file with \texttt{try}/\texttt{except FileNotFoundError} and explain why \texttt{"w"} can wipe data. & L2.G6\\
7.4 & The student can decide what should and should not be saved about a player, and keep personal information out of files. & L2.G9\\
\midrule
\multicolumn{3}{@{}l}{\textbf{8: Testing: Prove It Works}}\\
8.1 & The student can write a test function with \texttt{assert} that checks what a function returns. & L2.G6\\
8.2 & The student can run tests with pytest in Thonny and read \texttt{PASSED}, \texttt{FAILED} and the failure message. & L2.G6\\
8.3 & The student can choose normal and edge test cases (empty list, first, last, exactly on a limit) that catch hidden bugs. & L2.G3, L2.G6\\
8.4 & The student can use failing tests to find and fix bugs, and explain why people need to be able to trust software. & L2.G6, L2.G9\\
\midrule
\multicolumn{3}{@{}l}{\textbf{9: Event-Driven Games}}\\
9.1 & The student can connect keys to handler functions with \texttt{onkey} and \texttt{listen}, passing the function without parentheses. & L2.G7\\
9.2 & The student can build a game loop with \texttt{ontimer} that moves a list of objects one frame at a time and stops on game over. & L2.G2, L2.G7\\
9.3 & The student can detect a collision with a distance check and keep the game state in a dictionary. & L2.G3, L2.G4\\
9.4 & The student can choose a final-project option and describe its idea, and explain why borrowed images and sounds need a license and credit. & L2.G8, L2.G9\\
\midrule
\multicolumn{3}{@{}l}{\textbf{10: Project Studio 1}}\\
10.1 & The student can turn a project idea into pseudocode or a flowchart and point to the loop, the decisions and the data in it. & L2.G8\\
10.2 & The student can break the project into named functions and split the tasks fairly with a partner on a task board. & L2.G3, L2.G8\\
10.3 & The student can build a core loop that runs without crashing: a game timer loop or an art-studio menu loop (Milestone M2). & L2.G2, L2.G7\\
10.4 & The student can give kind, specific code-review feedback and act on feedback from a partner. & L2.G8\\
\midrule
\multicolumn{3}{@{}l}{\textbf{11: Project Studio 2}}\\
11.1 & The student can write at least 3 pytest tests (normal and edge cases) for the team's own functions and run them. & L2.G6\\
11.2 & The student can add a list, a dictionary and a file that saves and loads to the project (Milestone M3). & L2.G4, L2.G5\\
11.3 & The student can run a playtest, record feedback on cards and turn it into a bug list ordered by importance. & L2.G8\\
11.4 & The student can check the project for accessibility and make at least one improvement (Milestone M4). & L2.G9\\
\midrule
\multicolumn{3}{@{}l}{\textbf{12: Showcase and Celebration}}\\
12.1 & The student can present the team's project and explain one design decision and one piece of code (a function, the loop, the list, the dictionary or the file). & L2.G8\\
12.2 & The student can give kind, specific feedback to other makers during a gallery walk. & L2.G8\\
12.3 & The student can show what they learned in Level 2 on the exit assessment: files, tests, events and functions. & L2.G3, L2.G5, L2.G6, L2.G7\\
12.4 & The student can explain how to share a project safely: family only, nicknames, no personal information on screen. & L2.G9\\
\midrule
\multicolumn{3}{@{}l}{\textbf{B1: Cipher \& Puzzle Day}}\\
B1.1 & The student can encode and decode a message with a Caesar shift, a number code and a keyword cipher by following the steps exactly. & L2.G8\\
B1.2 & The student can crack a Caesar message without the key by trying shifts in order or by spotting the most common letter. & L2.G2\\
B1.3 & The student can solve a logic puzzle by writing down facts and ruling out possibilities one at a time. & L2.G8\\
B1.4 & The student can spot a message that tries to trick them into sharing a password or secret. & L2.G9\\
\midrule
\multicolumn{3}{@{}l}{\textbf{B2: Generative Art Jam}}\\
B2.1 & The student can use a list of colors and random choices to make a picture that changes every run. & L2.G4, L2.G7\\
B2.2 & The student can write a drawing function with parameters and call it many times from a loop. & L2.G3\\
B2.3 & The student can use \texttt{random.seed} to get the same \enquote{random} picture again. & L2.G1, L2.G7\\
B2.4 & The student can give credit when remixing someone else's art idea or code. & L2.G9\\
\midrule
\bottomrule\end{longtable}
